% Options for packages loaded elsewhere
\PassOptionsToPackage{unicode}{hyperref}
\PassOptionsToPackage{hyphens}{url}
\documentclass[
]{article}
\usepackage{xcolor}
\usepackage{amsmath,amssymb}
\setcounter{secnumdepth}{-\maxdimen} % remove section numbering
\usepackage{iftex}
\ifPDFTeX
  \usepackage[T1]{fontenc}
  \usepackage[utf8]{inputenc}
  \usepackage{textcomp} % provide euro and other symbols
\else % if luatex or xetex
  \usepackage{unicode-math} % this also loads fontspec
  \defaultfontfeatures{Scale=MatchLowercase}
  \defaultfontfeatures[\rmfamily]{Ligatures=TeX,Scale=1}
\fi
\usepackage{lmodern}
\ifPDFTeX\else
  % xetex/luatex font selection
\fi
% Use upquote if available, for straight quotes in verbatim environments
\IfFileExists{upquote.sty}{\usepackage{upquote}}{}
\IfFileExists{microtype.sty}{% use microtype if available
  \usepackage[]{microtype}
  \UseMicrotypeSet[protrusion]{basicmath} % disable protrusion for tt fonts
}{}
\makeatletter
\@ifundefined{KOMAClassName}{% if non-KOMA class
  \IfFileExists{parskip.sty}{%
    \usepackage{parskip}
  }{% else
    \setlength{\parindent}{0pt}
    \setlength{\parskip}{6pt plus 2pt minus 1pt}}
}{% if KOMA class
  \KOMAoptions{parskip=half}}
\makeatother
\setlength{\emergencystretch}{3em} % prevent overfull lines
\providecommand{\tightlist}{%
  \setlength{\itemsep}{0pt}\setlength{\parskip}{0pt}}
\usepackage{bookmark}
\IfFileExists{xurl.sty}{\usepackage{xurl}}{} % add URL line breaks if available
\urlstyle{same}
% fallback for those not using the hyperref driver hyperxmp:
\makeatletter
\@ifundefined{xmpquote}{\newcommand{\xmpquote}[1]{#1}}{}
\makeatother
\hypersetup{
  pdftitle={Python Programming},
  hidelinks,
  pdfcreator={LaTeX via pandoc}}

\title{\texttt{Python} Programming}
\author{}
\date{}

\begin{document}
\maketitle

\subsection{Description}\label{description}

\texttt{Python} is an easy-to-use interpreted language that has steadily
gained in popularity over the last few years in a wide spectrum of
applications, ranging from \texttt{AI} to \texttt{Web\ Services}.
\texttt{Python} is also powerful, portable, object-oriented open source
programming language for writing stand alone programs, quick scripts,
and prototypes for large applications. This course provides an in-depth
and hands-on introduction to the \texttt{Python} programming language,
as well as the most common \texttt{Python} application domains and
tools.

This course deals with \texttt{Python} version 3. \texttt{Python}
version 2 is not discussed.

\subsection{Duration}\label{duration}

40 hours / 5 days

\subsection{Intended Audience}\label{intended-audience}

\begin{itemize}
\tightlist
\item
  Developers who need a flexible and multi purpose scripting language at
  their disposal
\item
  System administrators who need the powerful scripting capabilities of
  \texttt{Python}.
\item
  \texttt{DevOps} people who need \texttt{Python} for their day to day
  operations.
\item
  Developers who plan to use \texttt{Python} as their primary
  programming language to build large scale systems.
\item
  Data scientists and \texttt{Machine\ Learning} experts who wish to
  enhance their toolkits with \texttt{Python}.
\end{itemize}

\subsection{Prerequisites}\label{prerequisites}

\begin{itemize}
\tightlist
\item
  Prior programming experience in another language such as Go,
  \texttt{Java}, \texttt{C\#} or C/\texttt{C++} is an advantage.
\end{itemize}

\subsection{Objectives}\label{objectives}

\begin{itemize}
\tightlist
\item
  Understand how \texttt{Python} works behind the scenes
\item
  Write short scripts for automation and testing
\item
  Get the basic skills to start writing larger \texttt{Python} systems
\item
  Could automate their work where-ever possible using \texttt{Python}
\item
  Could use \texttt{Python} for testing, web testing, building and
  deploying software
\item
  Will have the basic knowledge to start developing large systems using
  \texttt{Python}
\end{itemize}

\subsection{Exercises}\label{exercises}

\begin{itemize}
\tightlist
\item
  Done with \texttt{PyCharm} community on either \texttt{Windows} or
  \texttt{Linux} as the students need.
\end{itemize}

\subsection{Outline}\label{outline}

\begin{itemize}
\tightlist
\item
  Introduction to \texttt{Python} (1 hour)

  \begin{itemize}
  \tightlist
  \item
    History of \texttt{Python}
  \item
    Why \texttt{Python}?
  \item
    \texttt{Python} is an \texttt{OO} language
  \item
    \texttt{Python} is a dynamic language (what does that mean?)
  \item
    Basic \texttt{OO} principles (for those who need them)
  \item
    Comparison of \texttt{Python} with other programming languages in
    terms of performance, garbage collection, multi-\texttt{threading},
    learning curve and more
  \end{itemize}
\item
  Your first \texttt{Python} program (1 hour)

  \begin{itemize}
  \tightlist
  \item
    How to install the \texttt{Python} programming environment

    \begin{itemize}
    \tightlist
    \item
      Difference versions of \texttt{Python} and which to choose
    \item
      Sometimes you don't get to choose
    \end{itemize}
  \item
    Using \texttt{python} interactively
  \item
    Your first script
  \item
    Running your script
  \item
    \texttt{IDEs} and tools for \texttt{Python}
  \end{itemize}
\item
  Basic types and operators (2 hours)

  \begin{itemize}
  \tightlist
  \item
    Why do we need basic types

    \begin{itemize}
    \tightlist
    \item
      \texttt{int}
    \item
      float
    \end{itemize}
  \item
    \texttt{str}

    \begin{itemize}
    \tightlist
    \item
      Different quotes (\texttt{\textquotesingle{}}, \texttt{"},
      \texttt{"""})
    \item
      Concatenating strings
    \item
      String formatting and formatting aids
    \item
      String slices
    \item
      Splitting and joining * \texttt{Unicode} in \texttt{Python}
    \item
      \texttt{bool}
    \end{itemize}
  \end{itemize}
\item
  Basic Data structures (4 hours)

  \begin{itemize}
  \tightlist
  \item
    list

    \begin{itemize}
    \tightlist
    \item
      Literal syntax
    \item
      \texttt{OO} syntax
    \item
      Slices
    \item
      Catenating
    \item
      Deleting parts of lists
    \item
      Correct iteration
    \end{itemize}
  \item
    dict

    \begin{itemize}
    \tightlist
    \item
      Literal syntax
    \item
      \texttt{OO} syntax
    \item
      getting keys
    \item
      Deleting keys
    \item
      getting values
    \item
      Merging
    \item
      Correct iteration
    \end{itemize}
  \item
    set

    \begin{itemize}
    \tightlist
    \item
      Literal syntax
    \item
      \texttt{OO} syntax
    \item
      Merging
    \item
      Intersecting
    \item
      Subtracting
    \end{itemize}
  \item
    tuple

    \begin{itemize}
    \tightlist
    \item
      creating
    \item
      iterating
    \item
      tuples are immutable (sort of)
    \end{itemize}
  \item
    performances of these data structures
  \item
    Toolkit iteration

    \begin{itemize}
    \tightlist
    \item
      enumerate
    \item
      \texttt{zip}
    \end{itemize}
  \end{itemize}
\item
  Basic statements (2 hours)

  \begin{itemize}
  \tightlist
  \item
    Assignments
  \item
    Expressions
  \item
    print
  \item
    Conditionals if
  \item
    Loops (\texttt{while-loop}, \texttt{for-loop})
  \item
    range
  \item
    break and continue
  \item
    pass
  \end{itemize}
\item
  Functions (4 hours)

  \begin{itemize}
  \tightlist
  \item
    Why do we need functions
  \item
    defining functions with the def keyword
  \item
    default values for arguments
  \item
    arbitrary argument lists: *args
  \item
    keyword arguments: **kwargs
  \item
    \texttt{lambda} expressions
  \item
    Scoping: global vs function scopes
  \item
    Nested functions and closures
  \end{itemize}
\item
  Using Modules (2 hours)

  \begin{itemize}
  \tightlist
  \item
    Why do we need modules
  \item
    Namespaces
  \item
    The various ways of importing other modules
  \item
    Reloading modules
  \item
    Module search path
  \item
    Compiling modules
  \item
    The standard library of modules and it's documentation
  \item
    The dir function
  \end{itemize}
\item
  Classes (4 hours)

  \begin{itemize}
  \tightlist
  \item
    Why do we need classes
  \item
    \texttt{Python} Namespaces
  \item
    Object properties
  \item
    The class statement and class object

    \begin{itemize}
    \tightlist
    \item
      Constructing classes
    \end{itemize}
  \item
    Creating instances
  \item
    Writing methods
  \item
    Special methods \_\_ and operator overloading
  \item
    Instance objects
  \item
    Method objects
  \item
    Using class methods
  \item
    Inheritance in \texttt{Python}
  \item
    Mixing classes in \texttt{Python}
  \item
    Operator overloading
  \item
    Namespace lookup rules
  \item
    Design using classes
  \item
    Properties and decorators
  \end{itemize}
\item
  Exceptions (4 hours)

  \begin{itemize}
  \tightlist
  \item
    Why do we need exceptions
  \item
    How are exceptions used
  \item
    Handling exceptions
  \item
    Raising exceptions
  \item
    User defined exceptions
  \item
    Catching modes
  \item
    Cleanup actions with finally
  \end{itemize}
\item
  Functional \texttt{Python} (2 hours)

  \begin{itemize}
  \tightlist
  \item
    Iterators
  \item
    Generators
  \end{itemize}
\item
  Systems programming in \texttt{Python} (4 hours)

  \begin{itemize}
  \tightlist
  \item
    IO

    \begin{itemize}
    \tightlist
    \item
      formatted printing
    \item
      using files
    \item
      reading and writing text files
    \item
      reading and writing \texttt{JSON} or \texttt{YAML} files
    \end{itemize}
  \item
    Using the \texttt{subprocess} module
  \item
    Using the \texttt{multiprocessing} module
  \item
    Using the \texttt{threading} module
  \end{itemize}
\item
  Using third party modules and packaging modules (2 hours)

  \begin{itemize}
  \tightlist
  \item
    Built in \texttt{Python} modules

    \begin{itemize}
    \tightlist
    \item
      pip
    \end{itemize}
  \item
    Downloading, installing and using modules off the net
  \end{itemize}
\item
  Creating Modules (2 hours)

  \begin{itemize}
  \tightlist
  \item
    Directory structure
  \item
    \textbf{init}.py file and code in it
  \item
    Using nearby modules
  \item
    Global code in your module
  \item
    Writing functions and classes in your module
  \item
    Documenting your module using \texttt{pydoc}

    \begin{itemize}
    \tightlist
    \item
      Testing your modules using \texttt{pytest}
    \end{itemize}
  \item
    Distributing modules using \texttt{setuptools} and
    \texttt{distutils}
  \end{itemize}
\item
  Technology around \texttt{Python} (2 hours)

  \begin{itemize}
  \tightlist
  \item
    Using the \texttt{Python} debugger
  \item
    Using the \texttt{Python} profiler
  \item
    Using different repositories than \texttt{PyPI}
  \item
    Using \texttt{virtualenv} and other tools
  \end{itemize}
\end{itemize}

\subsection{Installations}\label{installations}

Each student should have:

\begin{itemize}
\tightlist
\item
  The most recommended is an \texttt{Ubuntu} 24:04 machine, real or
  virtual. The reason for this is that most companies production
  environment is \texttt{Linux} and not \texttt{Windows} or
  \texttt{MacOS} and so it is better to practice on an environment as
  similar as possible to the production environment as opposed to an
  environment used to read emails.
\item
  4 GB \texttt{RAM} for each machine is enough. This is not much.
\item
  Free, wide band, access to the internet from all machines with no
  weird corporate firewalls that might stop us from installing software
  and \texttt{Python} packages via pip.
\item
  Username and password of a user that has \texttt{sudo} privileges on
  the machine.
\item
  \url{https://www.linuxvmimages.com/images/ubuntu-2204/}
\item
  While \texttt{Ubuntu} is the best to exercise \texttt{Python} users
  who want to exercise on \texttt{Windows} or on \texttt{MacOS} are
  welcome to do so and I will help them out with issues.
\item
  No need to install \texttt{Python} yourself as I will help the
  students to install \texttt{Python} when the course starts.
\item
  Students who wish to practice on \texttt{Windows} may do so, but be
  sure that you can install \texttt{Python} on your machine (some
  companies and enterprises prohibit installation of software on
  developers machines). Be aware that \texttt{Python} was designed with
  \texttt{UNIX} in mind and so some systems programming features may not
  work correctly on \texttt{Windows} but impacts this beginners course
  very little.
\item
  Students who wish to practice on \texttt{MacOS} may do so as well.
  Again, \texttt{make} sure that you have \texttt{Python} installed on
  your laptop (usually it is already installed) or that you have the
  privilege to install it.
\item
  Recommended development environments: \texttt{vscode} with
  \texttt{Python} plugins, \texttt{PyCharm} community edition from
  \texttt{JetBrains}.
\end{itemize}

\subsection{Copyright}\label{copyright}

Mark Veltzer
\href{mailto:mark.veltzer@gmail.com}{\nolinkurl{mark.veltzer@gmail.com}},
© 2026

\end{document}
